\chapter{From \texorpdfstring{$+1.4$}{+1.4} to a probability: the currency of evaluation}
\label{ch:currency}

\begin{objectives}
\begin{itemize}
\item Say what a centipawn evaluation actually claims, and why the claim is empirical rather
  than definitional.
\item Build the logistic curve from scratch and use it to convert an evaluation into a
  win probability.
\item Read a WDL triple, convert between $(w,d,l)$, expected score, and $V$, and do it on
  real Leela output.
\item Explain --- with arithmetic --- why Monte-Carlo search requires probabilities and
  would be broken by centipawns.
\item Distinguish two positions with identical evaluations but opposite \emph{character},
  and quantify the difference.
\end{itemize}
\end{objectives}

\section{What does $+1.4$ mean?}

A classical engine says $+1.4$. Ask what it means and you get "White is better by about a
pawn and a half" --- which just restates the units. The honest answer is that a centipawn
score means nothing on its own; it acquires meaning only through a claim about outcomes:
\emph{positions that score $+1.4$ tend to be won by White at some particular rate.}

That is an empirical statement. Somebody has to go and measure it, and the answer depends on
who is playing --- $+1.4$ between two 1600s and $+1.4$ between two engines are different
propositions about the future. Centipawns are a \textbf{unit of account}, not a unit of
truth, and the exchange rate to reality has to be measured.

Leela cuts out the middle step. Its network does not estimate "how many pawns"; it estimates
\emph{the outcome distribution directly} --- how often this position is won, drawn, lost.
That single design choice is the reason its numbers can be reasoned about, averaged, and
combined, and it is what makes the search of Part~\textsc{ii} mathematically legitimate.

\section{The score of a game, and expected score}

Chess has three outcomes. Score them the way tournaments do:
\[
  \text{win} = 1, \qquad \text{draw} = \tfrac12, \qquad \text{loss} = 0 .
\]
If a position has win probability $w$, draw probability $d$ and loss probability $l$ (with
$w + d + l = 1$), then the \textbf{expected score} is the average points you collect:
\begin{equation}
  \label{eq:expscore}
  \E[\text{score}] \;=\; 1\cdot w \;+\; \tfrac12 \cdot d \;+\; 0 \cdot l
  \;=\; w + \tfrac{d}{2}.
\end{equation}

Engines prefer a symmetric scale running from $-1$ (lost) through $0$ (balanced) to $+1$
(won). Rescale by $V = 2\,\E[\text{score}] - 1$:
\[
V = 2\left(w + \tfrac{d}{2}\right) - 1 = 2w + d - 1 .
\]
Since $d = 1 - w - l$, this collapses to something you can compute in your head:

\begin{finalform}[The value of a position]
\begin{equation}
  \label{eq:VfromWDL}
  \boxed{\;V \;=\; w - l\;}
\end{equation}
Draw probability does not appear. $V$ is simply how much more likely you are to win than to
lose. It ranges over $[-1,+1]$ and is stated \emph{from the point of view of the side to
move} (Chapter~\ref{ch:problem}'s negamax convention).
\end{finalform}

\begin{worked}{reading real Leela output}
Leela's verbose output reports two of the three numbers: \texttt{WL} (which is exactly
$V = w-l$) and \texttt{D} (which is $d$). Recover $w$ and $l$ by solving the pair
$w - l = \texttt{WL}$, $w + l = 1 - \texttt{D}$:
\begin{equation}
  \label{eq:recoverwdl}
  w = \frac{1 - \texttt{D} + \texttt{WL}}{2}, \qquad
  l = \frac{1 - \texttt{D} - \texttt{WL}}{2}.
\end{equation}
For the initial position, the engine in \texttt{engine/} reports
$\texttt{WL} = +0.0946$, $\texttt{D} = 0.461$:
\[
w = \frac{1 - 0.461 + 0.0946}{2} = \frac{0.6336}{2} = 0.3168, \qquad
l = \frac{1 - 0.461 - 0.0946}{2} = \frac{0.4444}{2} = 0.2222 .
\]
\end{worked}

\begin{realdata}[title={Real engine output: five positions, four currencies}]
Produced by \texttt{tools/collect\_engine\_data.py} with the network
\texttt{engine/791556.pb.gz}; the Stockfish column is Stockfish 16.1 at depth 30, used as an
independent referee.

\begin{center}\small
\begin{tabular}{@{}lrrrrl@{}}
\toprule
Position & $w$ & $d$ & $l$ & $V = w-l$ & Stockfish d30 \\
\midrule
Initial position                 & 0.317 & 0.461 & 0.222 & $+0.095$ & $+0.35$ \\
Opposition draw (K+P, §2.6)      & 0.003 & 0.997 & 0.000 & $+0.003$ & $0.00$ \\
K+P, White to move (won)         & 0.942 & 0.058 & 0.000 & $+0.942$ & mate in 12 \\
Same, Black to move (lost)       & 0.000 & 0.052 & 0.948 & $-0.948$ & mate in 16 \\
Morphy position before 16.Qb8+   & 0.713 & 0.193 & 0.094 & $+0.620$ & \textbf{mate in 2} \\
\bottomrule
\end{tabular}
\end{center}
\end{realdata}

That last row is worth a pause, because it is a lesson the rest of the book keeps returning
to. In a position with a forced mate in two on the board, this network's evaluation is
$+0.62$ --- "clearly better, probably winning". It is not lying and it is not broken. It is
reporting \emph{how positions of this kind tend to go}, which is not the same as
\emph{what is true here}. Finding the mate is the search's job (you will watch it happen, in
real numbers, in Chapter~\ref{ch:readingtree}).

\section{Converting a centipawn score into a probability}

You will constantly need to move between the two currencies, since Stockfish speaks
centipawns and Leela speaks probabilities. The bridge is a curve we can build from scratch.

\subsection{Building the logistic function}

We need a function that takes an evaluation $x$ (any real number, positive or negative) and
returns a probability (between 0 and 1). Build it in three steps.

\begin{attempt}{1}{a straight line}
$p = \tfrac12 + kx$. Fails immediately: at $x = +10$ it returns a probability above 1. Any
straight line eventually leaves $[0,1]$. We need something that \emph{saturates}.
\end{attempt}

\begin{attempt}{2}{work with odds instead}
Instead of probability, use \textbf{odds} $= p/(1-p)$: a probability of $0.75$ is odds of
$3$. Odds live in $(0,\infty)$ --- half the problem gone. And odds \emph{multiply} naturally:
doubling your advantage should double your odds, not add a fixed amount of probability.
Take logs and the multiplication becomes addition:
\[
\log\text{-odds} = \ln\frac{p}{1-p} \in (-\infty, +\infty).
\]
Now \emph{that} is a scale on which a linear model is sensible. So posit the simplest
possible relationship:
\[
\ln\frac{p}{1-p} \;=\; k\,x .
\]
\end{attempt}

\begin{attempt}{3}{solve it back for $p$}
Exponentiate: $\frac{p}{1-p} = e^{kx}$, so $p = (1-p)e^{kx}$, so $p(1 + e^{kx}) = e^{kx}$:
\[
p = \frac{e^{kx}}{1+e^{kx}} = \frac{1}{1+e^{-kx}} .
\]
\end{attempt}

\begin{finalform}[The logistic (sigmoid) function]
\begin{equation}
  \label{eq:logistic}
  \sigma(z) \;=\; \frac{1}{1+e^{-z}}, \qquad
  p(\text{win-ish}) \;=\; \sigma(k\,x)
\end{equation}
with $x$ the evaluation and $k$ a scale constant \emph{fitted to data}, not derived. Note
$\sigma(0) = 1/2$, $\sigma(+\infty) = 1$, $\sigma(-\infty) = 0$, and the symmetry
$\sigma(-z) = 1 - \sigma(z)$.
\end{finalform}

You will meet $\sigma$ three more times in this book, wearing different hats: as the output
of a value head (Chapter~\ref{ch:heads}), as the thing that turns a linear probe's raw score
into a probability (Chapter~\ref{ch:probes}), and as the single-neuron ancestor of the whole
network (Chapter~\ref{ch:learned}). It is the same function every time. Learn it once.

\begin{worked}{calibrating and checking the conversion}
A common fitted value (used by Lichess for its evaluation bars) is
$\text{expected score} = \sigma(0.00368\,\text{cp})$, with cp in centipawns.

Take Stockfish's $+35$ centipawns for the initial position:
\[
\sigma(0.00368 \times 35) = \sigma(0.1288) = \frac{1}{1+e^{-0.1288}}
= \frac{1}{1+0.8792} = 0.5321 .
\]
Expected score $0.532$, so by Equation~\eqref{eq:expscore} rescaled,
$V = 2(0.532)-1 = +0.064$.

Leela's own answer for the same position was $V = +0.095$ (§2.2). Two engines, two
architectures, two definitions --- and they agree that the first move is worth about
$0.06$--$0.10$ of a game result, i.e.\ that White scores about $53\%$ from the start. That is
a genuine, checkable cross-validation of the conversion, and about as much precision as the
conversion deserves.
\end{worked}

\begin{pitfall}
The constant $k$ is not a law of nature. It depends on \emph{whose} games were used to fit
it: at club level, a $+1.5$ position is won far less reliably than between engines, so the
same centipawn number maps to a different probability. Any tool that displays "White wins
$78\%$" must be able to say \emph{$78\%$ of the time, between whom}. This is the first place
in this book where a number can be technically correct and pedagogically dishonest, and it
will not be the last.
\end{pitfall}

\section{Why the search needs probabilities}

Here is the argument that makes this chapter structural rather than cosmetic. Monte-Carlo
search does exactly one thing with the numbers coming back from the network: it
\textbf{averages} them (Chapter~\ref{ch:onemachine}, $Q = W/N$). So the question is whether
averaging is a meaningful operation in each currency.

\begin{worked}{averaging in two currencies}
A move leads to two equally likely continuations. In one, you emerge a queen up; in the
other, you are mated.

\textbf{In centipawns:} the two leaves score $+900$ and $-\infty$ (mate). What is the average?
The question is not even well-posed --- engines have to fake it with a large finite number
like $\pm 32000$, and then the average of $+900$ and $-32000$ is $-15550$, a number that
corresponds to no position and no fact. Worse, the scale is not linear in anything: being
$+900$ is not "nine times better" than being $+100$, because both are usually just winning.

\textbf{In probabilities:} the two leaves score $V = +1$ and $V = -1$. The average is $0$,
and the average \emph{means something exact}: half the futures examined are won and half are
lost, so the expected score is $\tfrac12$ --- a coin flip. That is a true and useful summary
of the position, and it is the number a player would want.
\end{worked}

\begin{keyidea}
Probabilities are \emph{averageable}: the average of a set of win probabilities is itself a
win probability, and it is the win probability you would get by picking one of those futures
at random. Centipawns have no such property. Since Monte-Carlo search \emph{is} averaging,
Leela's value head is not a stylistic choice --- it is the thing that makes the search
coherent. An engine cannot back up centipawns through a Monte-Carlo tree without inventing
meaning it does not have.
\end{keyidea}

There is a caveat that costs a paragraph, because it will come back in
Chapter~\ref{ch:engineering}. Averaging is the right operation for "what does a random
future look like", but it is the \emph{wrong} operation for "what happens with best play".
If one of my two continuations is a forced mate against me, the position is lost, not a coin
flip --- my opponent will simply choose it. The average of $+1$ and $-1$ being $0$ is exactly
the minimax-versus-averaging tension, and every Monte-Carlo engine has to manage it. It is
managed, not solved: the search visits the losing branch more and more, dragging the average
down, so the answer converges to the truth without ever being computed as a minimum. The
speed of that convergence is a real limitation of the method and one of the honest
disadvantages of Leela versus alpha--beta on sharp forcing positions.

\section{Sharpness: two positions, one evaluation}

Because $V = w - l$ discards $d$, a single number cannot tell you what kind of fight you are
in. Two positions can share $V \approx 0$ and be nothing alike:

\begin{center}
\begin{tabular}{@{}lrrrrl@{}}
\toprule
 & $w$ & $d$ & $l$ & $V$ & Character \\
\midrule
Position A & 0.05 & 0.90 & 0.05 & $0.00$ & dead level; hard to lose, hard to win \\
Position B & 0.45 & 0.10 & 0.45 & $0.00$ & mutual minefield; someone is going to be wrong \\
\bottomrule
\end{tabular}
\end{center}

\emph{(Illustrative numbers.)} These demand opposite treatment from a training tool, and
opposite treatment from you at the board. A single evaluation bar cannot distinguish them.
The WDL triple can, and this is the cheapest, most under-used signal Leela produces.

Define a \textbf{sharpness} measure as the probability that the game is \emph{decided}:
\begin{equation}
  \label{eq:sharpness}
  \text{sharp}(s) \;=\; w + l \;=\; 1 - d .
\end{equation}
Position A scores $0.10$; position B scores $0.90$. Same $V$, ninefold difference in what is
at stake. And it is not a hand-waving construction --- the engine really does report it:

\begin{realdata}[title={Real engine output: draw probability separates the positions}]
\begin{center}
\begin{tabular}{@{}lrrl@{}}
\toprule
Position & $d$ & sharp $= 1-d$ & \\
\midrule
Opposition draw (K+P)           & 0.997 & 0.003 & nothing can happen \\
Initial position                & 0.461 & 0.539 & a normal game \\
Morphy position (before Qb8+)   & 0.193 & 0.807 & something is about to be decided \\
K+P, White to move              & 0.058 & 0.942 & already decided \\
\bottomrule
\end{tabular}
\end{center}
Note that the last two rows have very different $V$ ($+0.62$ vs $+0.94$) but similar
sharpness --- and the first two have similar $V$ (both near $0$) and utterly different
sharpness. The two numbers are genuinely independent axes.
\end{realdata}

\begin{keyidea}
$V$ tells you \emph{who} stands better. $1-d$ tells you \emph{whether it matters yet}. Your
project's interest in "sharp, dynamic chess" is, in this vocabulary, an interest in
positions with high $1-d$ --- and Chapter~\ref{ch:style} shows that you can steer an engine
towards them by writing that preference into an objective function.
\end{keyidea}

\subsection{What Leela's other columns mean}

While we are here, the rest of a verbose line, so nothing in this book is mysterious later:

\begin{center}
\begin{tabular}{@{}ll@{}}
\toprule
\texttt{V} & the network's raw evaluation of this move's position: one forward pass, no search \\
\texttt{WL} & $w - l$ of the current node, i.e.\ $V$ after search \\
\texttt{D} & draw probability $d$ \\
\texttt{M} & predicted number of plies remaining in the game (a third head; see ch.~\ref{ch:heads}) \\
\texttt{N} & visits (Chapter~\ref{ch:onemachine}) \\
\texttt{P} & policy prior (Chapter~\ref{ch:puct}) \\
\texttt{Q} & running average of everything backed up through this move \\
\texttt{U} & the exploration bonus (Chapter~\ref{ch:puct}) \\
\texttt{S} & $Q + U$: the PUCT score that actually decides where to look next \\
\bottomrule
\end{tabular}
\end{center}

By Chapter~\ref{ch:byhand} you will be able to compute \texttt{U} and \texttt{S} yourself
from the other columns, and check the engine's arithmetic. You will find it checks out to
five decimal places.

\begin{repobox}[title={in this repository}]
The WDL triple is already plumbed through this project: the frontend's evaluation display and
the "character of the fight" language in the design documents both come from $d$. The
sharpness quantity $1-d$ of Equation~\eqref{eq:sharpness} is the simplest ancestor of the
tactical-steering score you build in Chapter~\ref{ch:sharpness}.
\end{repobox}

\begin{exercises}

\exhead[conversions]{ch2:convert} A position has $w = 0.55$, $d = 0.30$, $l = 0.15$. Compute
the expected score and $V$. Another has $w = 0.40$, $d = 0.60$, $l = 0.00$. Which has the
larger $V$? Which would you rather have, and why might the answer depend on the tournament
situation rather than the position?

\exhead[recovering WDL]{ch2:recover} Leela reports $\texttt{WL} = -0.312$, $\texttt{D} =
0.244$. Recover $w$, $d$, $l$ with Equation~\eqref{eq:recoverwdl} and check they sum to 1.
Whose point of view are these from, and what changes if it is Black to move?

\exhead[the logistic by hand]{ch2:logistic} Compute $\sigma(z)$ for $z = -2, -1, 0, 1, 2$.
Sketch the curve. Where is it steepest, and what does that steepness say about which
positions an extra tenth of a pawn matters most in?

\exhead[centipawns to probability]{ch2:cp2p} Using $\sigma(0.00368\,\text{cp})$, convert
$+50$, $+150$, $+300$ and $+800$ centipawns into expected scores. Compute the difference in
expected score between $+50$ and $+150$, and between $+700$ and $+800$. Both are a
100-centipawn gain --- explain in one sentence why they are worth wildly different amounts.

\exhead[the ratings problem]{ch2:ratings} $k$ in Equation~\eqref{eq:logistic} is fitted to a
population. Should $k$ be larger or smaller when fitting to 1200-rated players than to
engines? Draw both curves on the same sketch, and say what this implies for a training tool
that shows you a win probability.

\exhead[averaging]{ch2:averaging} A move has four sampled continuations with
$V = +0.9, +0.8, +0.85, -1.0$. Compute the average. Now suppose the $-1.0$ is a forced mate
against you that the opponent can choose. What is the position actually worth? Explain the
discrepancy in terms of §2.4, and say what the search must do to fix it.

\exhead[sharpness]{ch2:sharp} Two positions: $(w,d,l) = (0.30,0.55,0.15)$ and
$(0.30,0.20,0.50)$. Compute $V$ and $1-d$ for both. Write the one-sentence description a
coach should give of each, using no chess terms at all.

\exhead[designing a display]{ch2:display} You must show a position's evaluation in your
tool using at most two visual elements. Argue for a specific choice of the two quantities to
display, given a user who wants sharp, dynamic play. What would be lost by showing the
classical single bar?

\exhead[the mate that reads as $+0.62$]{ch2:mate62} In §2.2, a position with mate in two
was evaluated by the network at $V = +0.62$. Explain why this is not a bug, using the word
"population" in your answer. Then say what would have to change for the number to become
$+1.0$, and which chapter of this book is about that change.

\exhead[a calibration experiment]{ch2:calibration} Design (do not run) an experiment on your
9{,}030-game corpus that would measure the true $k$ for \emph{your} games: what would you
compute, what would you plot, and what would the plot look like if the standard $k$ were too
large? Be specific about the buckets.

\end{exercises}
